\documentclass[11pt,a4paper]{article}
\usepackage{iftex}\ifPDFTeX\usepackage[T1]{fontenc}\fi
\usepackage{amsmath,amssymb,amsthm}
\usepackage[vmargin=29mm,hmargin=28mm]{geometry}
\usepackage[hidelinks]{hyperref}
\usepackage{url}
\newtheorem{theorem}{Theorem}[section]
\newtheorem{lemma}[theorem]{Lemma}
\newcommand{\dd}{\,\mathrm d}
\newcommand{\ii}{\mathrm i}
\newcommand{\erfc}{\operatorname{erfc}}
\newcommand{\e}{\mathrm e}
\setlength{\emergencystretch}{2em}
\title{Explicit Uniform Remainder Estimate for the Alternating Hurwitz Zeta Function in a Joint Large-Parameter Regime}
\author{Tiago Manuel Serafim Campina da Silva}
\date{\small Version date: 24 November 2026}
% Keep section counters for theorem numbers and hyperlink destinations.
\makeatletter
\renewcommand\@seccntformat[1]{}
\renewcommand\section{\@startsection{section}{1}{\z@}
  {-2.3ex plus -.4ex minus -.2ex}{1.1ex plus .2ex}
  {\normalfont\large\bfseries\centering}}
\renewcommand\subsection{\@startsection{subsection}{2}{\z@}
  {-2ex plus -.3ex minus -.2ex}{.8ex plus .2ex}
  {\normalfont\normalsize\bfseries\centering}}
\renewcommand\subsubsection{\@startsection{subsubsection}{3}{\z@}
  {-1.8ex plus -.3ex minus -.2ex}{.6ex plus .2ex}
  {\normalfont\normalsize\itshape\centering}}
\makeatother
\raggedbottom
\widowpenalty=6000
\clubpenalty=6000
\hypersetup{pdftitle={Explicit Uniform Remainder Estimate for the Alternating Hurwitz Zeta Function in a Joint Large-Parameter Regime},pdfauthor={Tiago Manuel Serafim Campina da Silva},pdfdisplaydoctitle=true}
\DeclareRobustCommand{\suppfile}[2]{\href{https://github.com/tiagocampinasilva/alternating-hurwitz-zeta/blob/v1.0.0/#1}{#2}}
\begin{document}
\maketitle
\pdfbookmark[1]{Abstract}{abstract}
\begin{abstract}
We consider the alternating Hurwitz zeta function with the normalization
\[
 Z_q(x)=q^{\ii t}\zeta_E\!\left(\tfrac12+\ii t,q\right),
 \qquad t=\pi q+x\sqrt q .
\]
The leading term comes from a Fresnel integral and can be written in terms of the
complementary error function. The proposed two-term approximation includes a correction of order
$q^{-1/2}$ and has absolute error strictly below
$4.3/q$ for all real $q\ge100$ and $|x|\le1$.
The alternating Abel--Plana formula separates the integral giving the
transition profile from a residual kernel. Rotation through $-\pi/4$
gives an integrable bound for the Taylor remainder. For the residual
kernel we give both a bound obtained by splitting the integral at
$\log q$ and a bound valid on the whole positive axis.
\end{abstract}

\section{Introduction}
At height $t=\pi q$, the complex power in the alternating summation
formula cancels the exponential decay of its kernel to first order.
A change of order $\sqrt q$ in $t$ leaves a linear term in the scaled
exponent. Its coefficient, $x=(t-\pi q)/\sqrt q$, determines the
transition profile. We seek an explicit error bound for the profile
and its first correction over $-1\le x\le1$, starting at $q=100$.

The alternating series is the Lerch zeta function at phase $1/2$
\cite{Apostol1951,DLMF2026}. For our complex $s$, Hu and Kim \cite{HuKim2024}
give an integral for their truncation error (Theorem~3.5) and bounds on
its modulus (Theorem~3.7). At fixed truncation order, the right-hand
side of their bound (3.5) grows linearly with the shift along our
transition curve. Their Theorem~3.6 assumes real $s$.
For fixed $s$, the large-shift expansion of Ferreira and López
\cite{FerreiraLopez2004} includes the alternating case $z=-1$.
Cai and López \cite{CaiLopez2019} study the positive real cut. Their
large-shift remainder estimate is for $z>1$. For the Hurwitz zeta function,
Fernandez and Fokas \cite{FernandezFokas2018} derive an expansion to all
orders, and Nemes \cite{Nemes2017} bounds the large-shift remainder.

Balanzario \cite{Balanzario2004} estimates the error in a Riemann--Siegel
formula for the Hurwitz zeta function. Balanzario and Sánchez-Ortiz
\cite{BalanzarioSanchez2012} give an integral formula for the Lerch zeta
function suited to numerical quadrature. Approximate functional equations
offer another comparison. Garunkštis, Laurinčikas and Steuding
\cite{GarunkstisLaurincikasSteuding2003} treat the Lerch case, and Miyagawa
\cite{Miyagawa2017} treats both Hurwitz and Lerch. A comparison of constants
must match the phase and normalization and account for truncation and any
residues.

For the origins of the Riemann--Siegel formula see Riemann and Siegel
\cite{Riemann1859,Siegel1932}, and for later accounts see Titchmarsh
and Edwards \cite{Titchmarsh1986,Edwards1974}. Guthmann treats series
associated with cusp forms \cite{GuthmannCusp1997} and the square of the
Riemann zeta function \cite{GuthmannSquare1997}. The contour analysis
is related to classical asymptotics \cite{Olver1997}, particularly when a
stationary point approaches a singularity \cite{Bleistein1966} or an
endpoint \cite{FernandezSpenceFokas2018}. Paris studies Stokes phenomena
for the Lerch and Hurwitz functions \cite{Paris2016,ParisHurwitz2005}.

Balanzario \cite{Balanzario2006} gives an Euler--Maclaurin formula for the
Hurwitz function, and Butzer and coauthors \cite{ButzerEtAl2011} relate
Euler--Maclaurin to Abel--Plana summation. For arbitrary-precision evaluation
of the Lerch transcendent, Navas-Palencia \cite{NavasPalencia2023} uses
Euler--Maclaurin summation and uniform expansions.
\section{Statement and notation}
Throughout, $q$ and $x$ are real. Set
\begin{equation}\label{eq:parameters}
 h=q^{-1/2},\qquad t=\pi q+x\sqrt q,\qquad s=\tfrac12+\ii t .
\end{equation}
All complex powers use the principal logarithm. For $q>0$ and
$\Re s>0$ write
\begin{equation}\label{eq:zeta}
 \zeta_E(s,q)=\sum_{n=0}^\infty\frac{(-1)^n}{(n+q)^s}
 =2^{-s}\bigl[\zeta(s,q/2)-\zeta(s,(q+1)/2)\bigr].
\end{equation}
Pairing consecutive terms gives a locally uniformly convergent series
on $\Re s>0$, since each difference is bounded on compact subsets
by a constant times $(2n+q)^{-1-\Re s}$. The Hurwitz identity follows
by separating parity on $\Re s>1$ and continuing across the removable
pole of the difference. Define
\begin{align}
 Z_q(x)&=q^{\ii t}\zeta_E(s,q),\label{eq:Z}\\
 F(x)&=\frac{1+\ii}{2}\exp\!\left(-\frac{\ii x^2}{2\pi}\right)
 \erfc\!\left(-\frac{\e^{-\ii\pi/4}x}{\sqrt{2\pi}}\right),\label{eq:F}\\
 G(x)&=\frac12+\frac{\ii}{6\pi}
 -\frac{x^2}{6\pi^2}\bigl(1-\ii xF(x)\bigr).\label{eq:G}
\end{align}
\begin{theorem}[Proposed uniform estimate]\label{thm:main}
For $q\ge100$ and $|x|\le1$,
\begin{equation}\label{eq:main}
 \left|Z_q(x)-F(x)-\frac{G(x)}{\sqrt q}\right|<\frac{4.3}{q}.
\end{equation}
\end{theorem}
Write $R_q(x)$ for the difference inside the absolute value. The bound
also covers nonintegral $q$.

\section{The alternating Abel--Plana decomposition}
\begin{lemma}\label{lem:ap}
For real $q>0$ and $\Re s>0$,
\begin{equation}\label{eq:ap}
\zeta_E(s,q)=\tfrac12q^{-s}
+\ii\int_0^\infty\frac{(q+\ii y)^{-s}-(q-\ii y)^{-s}}
 {2\sinh(\pi y)}\dd y .
\end{equation}
\end{lemma}
\begin{proof}
Apply the Gamma integral to each term of a finite alternating sum
and pass to the limit to obtain
\[
 \Gamma(s)\zeta_E(s,q)=
 \int_0^\infty v^{s-1}\e^{-qv}(1+\e^{-v})^{-1}\dd v.
\]
The partial geometric sums have modulus at most two. Twice the modulus of the Gamma weight is integrable and bounds each finite-sum integrand. For $v>0$,
\[
 \int_0^\infty\frac{\sin(vy)}{\sinh(\pi y)}\dd y
 =\tfrac12\tanh(v/2).
\]
To evaluate the sine integral, integrate $\e^{ivz}/\sinh(\pi z)$ over a rectangle of height
$\ii$, indenting about $0$ and $\ii$ into its interior. The two
clockwise semicircles contribute $-\ii$ and $\ii\e^{-v}$.
The vertical sides vanish at infinity. If $I$ is the principal value
on the real axis, the horizontal integrals give $(1+\e^{-v})I$.
The imaginary part of the integrand is even, which gives the half-line formula. Now $(1+\e^{-v})^{-1}=(1+\tanh(v/2))/2$.
The bound $|\sin(vy)|\le vy$ permits interchange of integrals,
since $v^{\Re s}\e^{-qv}$ and $y/\sinh(\pi y)$ are integrable.
For fixed $y$ evaluate the two Gamma integrals from the sine, taking
their difference before integration in $y$. This gives \eqref{eq:ap}.
\end{proof}

Write $\theta=\arctan(y/q)$ and set
\[
 w=\e^{-\pi y}(1+\ii y/q)^{-s},\qquad
 v_-=\e^{-\pi y}(1-\ii y/q)^{-s},\qquad
 D=2\pi y-2t\theta+\ii\theta .
\]
Since $v_-/w=\exp(2\ii s\theta)$, define for $y>0$
\begin{equation}\label{eq:K}
 K=\frac{\e^{-2\pi y}w-v_-}{1-\e^{-2\pi y}}
   =\frac{w(1-\e^D)}{\e^{2\pi y}-1}.
\end{equation}
At the origin the quotient extends continuously with value $K(0)=-D'(0)/(2\pi)=x/(\pi\sqrt q)-\ii/(2\pi q)$.
Splitting $(w-v_-)/(1-\e^{-2\pi y})=w+K$ in \eqref{eq:ap} yields
\begin{equation}\label{eq:split}
 Z_q=h/2+\ii H(h,x)+\ii h\int_0^\infty K(y)\dd y,\qquad
 H(h,x)=h\int_0^\infty w(y)\dd y .
\end{equation}

\section{The residual-kernel bound}
From now on we assume the hypotheses of Theorem~\ref{thm:main}. In particular, $h\le1/10$ and $t>0$.
Writing $r=y/q$ and using $0\le r-\arctan r\le r^3/3$ gives
\begin{equation}\label{eq:D}
 |D|\le (2h+q^{-1})y+\frac{2\pi y^3}{3q^2},\qquad
 |w|\le\e^{hy},\qquad |v_-|\le\e^{-\pi y}.
\end{equation}
The imaginary part of $D$ is $\theta\le y/q$, and
$\pi q(r-\theta)$ and $|x|\sqrt q\,\theta$ bound the two contributions to
its real part after division by two. The moduli of the complex powers
contain the factors $\exp(\pm t\theta)(1+r^2)^{-1/4}$.

\begin{lemma}[Splitting at $\log q$]\label{lem:local}
The scaled residual integral is bounded by
$\sqrt q\int_0^\infty |K(y)|\dd y < B_{\rm s}=\frac{689147}{1575000}$.
\end{lemma}
\begin{proof}
On $0\le y\le L=\log q$ the bound for the phase difference in \eqref{eq:D}
increases with the integration variable, so it suffices to set $y=L$. Each term $q^{-a}\log^kq$ decreases for
$(a,k)=(1/2,1),(1,1),(2,3)$, since $k-a\log q<0$. Using
$\log100<4.61$ and $\pi<22/7$ gives
$|D|<51902943991/52500000000<0.989$ and $|w|\le\e^{0.461}<5/3$.
Thus
\[
 |w(1-\e^D)|\le (5/3)(\e-1)|D|<3|D|\le5|D|\qquad(D\ne0).
\]
At $D=0$ both sides vanish. We use the multiplier $5$ in the principal
estimate. On $y\ge L$, the first form of \eqref{eq:K}
and \eqref{eq:D} give $|K|\le3\e^{-\pi y}$, because
$2\pi-h>\pi$ and $1-\e^{-2\pi L}>2/3$.
Expanding $(\e^{2\pi y}-1)^{-1}$ into positive exponentials gives
the moments of orders one and three as $1/24$ and $1/240$.
In general the moment of order $j$ is $j!\zeta(j+1)/(2\pi)^{j+1}$.
Extending the short-interval majorants to infinity now gives
\[
 \sqrt q\int_0^\infty |K|\dd y\le
 \frac5{12}+\frac5{24\sqrt q}+\frac{5\pi}{360q^{3/2}}
 +\frac3\pi q^{1/2-\pi}.
\]
Since the exponent is smaller than minus five halves and the prefactor is less than one, the last term is less than $10^{-5}$. Bounding the preceding terms by
$5/12+5/240+22/(7\cdot72000)$ proves the lemma.
\end{proof}

\begin{lemma}[Global residual estimate]\label{lem:global}
For every $y>0$,
\begin{equation}\label{eq:whole}
 |K(y)|\le
 \frac{(2q^{-1/2}+q^{-1})y+2\pi y^3/(3q^2)}
 {2\sinh(\pi y)} .
\end{equation}
Consequently $\sqrt q\int_0^\infty |K|\dd y < B_{\rm g}=\frac{22061}{84000}$.
\end{lemma}
\begin{proof}
Integrating the derivative of $\exp(vD)$ over $0\le v\le1$ gives $|w(1-\e^D)|\le |D|\max(|w|,|w\e^D|)$. Each modulus in the maximum is at most $\e^{\pi y}$, the second because
$w\e^D=\e^{2\pi y}v_-$. Substitution into \eqref{eq:K} gives
\eqref{eq:whole}. At zero the bound has a finite limit, and at infinity
it decays exponentially. The positive exponential series for
$(2\sinh\pi y)^{-1}$ contains precisely the odd indices. Its moment
of order $j$ is $j!(1-2^{-j-1})\zeta(j+1)/\pi^{j+1}$, giving
$1/8$ and $1/16$ for $j=1,3$. Hence
\[
 \sqrt q\int_0^\infty |K|\dd y
 \le \frac{2+q^{-1/2}}8+\frac{\pi}{24q^{3/2}}
 <\frac{21}{80}+\frac{22}{7\cdot24000}=B_{\rm g}.
\]
\end{proof}

\section{Rotation of the principal integral}
Let $\Gamma=\e^{-\ii\pi/4}[0,\infty)$, oriented outwards. In the sector
$-\pi/4\le\arg y\le0$ the base $1+\ii y/q$ has real part at least one.
Close the contour with the triangle whose vertices are $0,R,R-\ii R$.
On its vertical side the integral of the modulus of $w$ is bounded by $R\exp(-\pi R+t\pi/2)$.
For fixed $q,x$ this tends to zero. On each ray the exponential in the real
part of the integration variable also gives absolute convergence. Cauchy's theorem and $z=hy$ give
\begin{equation}\label{eq:phase}
 H=\int_\Gamma \e^{\Psi_h(z)}\dd z,\quad
 \Psi_h=-\frac{\pi z}{h}
 -\left(\frac12+\frac{\ii\pi}{h^2}+\frac{\ii x}{h}\right)\log(1+\ii hz).
\end{equation}

\begin{lemma}[Regularity at $h=0$]\label{lem:tay}
On the ray $z=\e^{-\ii\pi/4}u$, $u\ge0$, the phase and its first two derivatives in real $h$ extend continuously
to $0\le h\le1/10$. Its constant and linear Taylor coefficients are
\[
 \Phi_x=xz-\ii\pi z^2/2,\qquad
 P_x=-\ii z/2-\ii xz^2/2-\pi z^3/3 .
\]
Moreover, $|\partial_h\Psi_h|\le u/2+u^2/2+\pi u^3/3$ and
$|\partial_h^2\Psi_h|\le u^2/2+2u^3/3+\pi u^4/2$.
\end{lemma}
\begin{proof}
Set $b=\ii z$, $L_m=\int_0^1 v^{m-1}/(1+hbv)\dd v$.
The exact identities $\log(1+hb)/h=bL_1$ and
$(\log(1+hb)-hb)/h^2=-b^2L_2$ turn \eqref{eq:phase} into
\[
 -\tfrac12\log(1+hb)-\ii xbL_1+\ii\pi b^2L_2 .
\]
The real part of $b$ is nonnegative.
For $n=0,1,2$, differentiation under the finite integral gives
$|L_m^{(n)}|\le n!u^n/(m+n)$ since $|1+hbv|\ge1$. Using also
$L_m(0)=1/m$, $L_m'(0)=-b/(m+1)$ gives the two coefficients displayed above.
\end{proof}

\begin{lemma}[Uniform integrable majorant]\label{lem:uniform}
Define
\begin{align}
 M(u)&=\tfrac74\e^{-9u^2/8}+\e^{-67u/20},\\
 P_b(u)&=\frac{484 u^6}{441}+\frac{22 u^5}{21}+\frac{241 u^4}{84}+\frac{7 u^3}{6}+\frac{3 u^2}{4} .
\end{align}
For the entire parameter window
$|\partial_h^2 \e^{\Psi_h}|\le M(u)P_b(u)$.
\end{lemma}
\begin{proof}
For $r=hu$, put $\vartheta(r)=\arctan(r/(\sqrt2+r))$ and
$\delta(r)=r/\sqrt2-\vartheta(r)$. Taking real parts in \eqref{eq:phase}
gives $\Re\Psi_h\le-\pi\delta(r)/h^2+u/\sqrt2$ for $h>0$.
Here $-\log|1+hb|/2\le0$ and $x\vartheta(r)/h\le u/\sqrt2$.
As $\vartheta'=1/[\sqrt2(1+\sqrt2r+r^2)]$ decreases, $\delta/r$
increases. Also
\[
 \frac{\delta'(r)}r=\frac{1+r/\sqrt2}{1+\sqrt2r+r^2}
\]
is decreasing. The denominator of its derivative is a positive square, and the numerator is $-1/\sqrt2-2r-r^2/\sqrt2$.
Integration shows that $\delta(r)/r^2$ decreases as well.

For an explicit transition point take $r_0=3/10$,
$\ell=140/99<\sqrt2<99/70=m$ and $v_0=r_0/(\ell+r_0)$.
The alternating series for the arctangent gives
$\delta(r_0)>d=r_0/m-v_0+v_0^3/3-v_0^5/5$.
Clearing the positive denominators gives the two comparisons needed below.
\[
 (157/50)d/r_0^2>9/8+5/22,\qquad
 10(157/50)d/r_0-1/\ell>67/20.
\]
For $r\le r_0$, monotonicity and $\pi>157/50$ give
$\Re\Psi_h\le-(9/8+5/22)u^2+u/\sqrt2\le-9u^2/8+11/20$.
For $r\ge r_0$, monotonicity of $\delta/r$ and $1/h\ge10$ give
$\Re\Psi_h\le-67u/20$. The first branch uses $\e^{11/20}<7/4$.
To check the exponential bounds for
$0<z\le1$, sum the series through degree six and bound its tail by $(z^7/7!)/(1-z/8)$. The inequality $\log100<4.61$ used earlier
follows from the exponential series through degree sixteen at $4.61$.

Thus $|\e^{\Psi_h}|\le M(u)$, including $h=0$ by continuity.
The chain rule gives a further factor
$|\partial_h\Psi_h|^2+|\partial_h^2\Psi_h|$. Squaring and adding
the bounds of Lemma~\ref{lem:tay}, with $\pi<22/7$, yields $P_b$.
A polynomial times either summand of $M$ is integrable.
\end{proof}

\section{The approximation and the final constant}
Set $J_0=\int_\Gamma\e^{\Phi_x}\dd z$,
$J_1=\int_\Gamma P_x\e^{\Phi_x}\dd z$, and
$J_*=\int_0^\infty M(u)P_b(u)\dd u$. Taylor's theorem under the
integral gives
\begin{equation}\label{eq:remainder}
 |H-J_0-hJ_1|\le h^2 J_*/2 .
\end{equation}
The absolute value of the integrand in the Taylor remainder is bounded by
$(h-v)M(u)P_b(u)$ on $0\le v\le h$ and $0\le u<\infty$.

On the rotated ray $J_0$ is a Gaussian integral with linear term
$x\e^{-\ii\pi/4}u$. Complete the square for a real linear coefficient. Analytic continuation
in that coefficient then gives $\ii J_0=F$.
Integration of $\partial_z\e^{\Phi_x}$ gives the boundary term $-1$,
hence $\pi F'=1-\ii xF$. Differentiation in $x$ is justified on compact
sets by Gaussian decay. Consequently $\pi F''=-\ii F-\ii xF'$ and
$\pi F'''=-2\ii F'-\ii xF''$. Inserting these identities into the
integral of $P_x$ gives
\begin{align*}
 \ii J_1&=-\ii F'/2-\ii xF''/2-\pi F'''/3\\
 &=\ii F'/6-\ii xF''/6
 =\ii/(6\pi)-x^2(1-\ii xF)/(6\pi^2).
\end{align*}
This identifies $G=1/2+\ii J_1$ in \eqref{eq:G}.

\begin{proof}[Proof of Theorem~\ref{thm:main}]
For $a=9/8$ and $\kappa=67/20$ let $g_j$ and $e_j$ be the moments of
$u^j$ against $\exp(-au^2)$ and $\exp(-\kappa u)$ on $[0,\infty)$.
Integration by parts gives
\[
 g_0=\tfrac12\sqrt{\pi/a}<67/80,\quad
 g_1=1/(2a),\quad g_j=(j-1)g_{j-2}/(2a)\ (j\ge2),
 \qquad e_j=j!/\kappa^{j+1}.
\]
If $p_j$ is the coefficient of $u^j$ in $P_b$, then
\[
 J_*=\sum_{j=2}^6 p_j(7g_j/4+e_j)
 \le \overline J=\frac{533948559110614545113}{69278297361197722560} .
\]
By \eqref{eq:split} and \eqref{eq:remainder}, the triangle inequality
gives $q|R_q(x)|\le J_*/2+\sqrt q\int_0^\infty|K|\dd y$. Using
Lemma~\ref{lem:local}, we obtain
\begin{equation}\label{eq:classicC}
 q|R_q(x)|\le\overline J/2+B_{\rm s}
 =\frac{371609064358084350147257}{86597871701497153200000}<43/10.
\end{equation}
Division by $q>0$ gives \eqref{eq:main}. The argument includes
$q=100$ and both endpoints $x=\pm1$.
\end{proof}

Substituting the whole-axis estimate of Lemma~\ref{lem:global} instead gives
\begin{equation}\label{eq:otherC}
 \overline J/2+B_{\rm g}=\frac{14258445238530472531037}{3463914868059886128000}=4.11628050389\ldots .
\end{equation}
Either residual estimate suffices for Theorem~\ref{thm:main}.

\section{Numerical cross-checks}\label{sec:numerics}
We evaluated \eqref{eq:zeta}--\eqref{eq:G} in \suppfile{Verification/verification.wl}{Wolfram Mathematica 15.0.1}
\cite{WolframMathematica} at the seven rational $q$ values in
Table~\ref{tab:max} and the \suppfile{Verification/verification_results.csv}{41 points} $x=-1+j/20$, $j=0,\ldots,40$.
We used 55 working digits and reevaluated the maximizing point in each row at 75 digits.
The results agree to the digits shown. Figure~\ref{fig:slice} shows the
slice with the largest sampled error.

The main error bound was evaluated using \suppfile{Verification/exact_coefficients.csv}{exact rational arithmetic}.
Mathematica was also used to verify the majorant polynomial and the
differential identity for the leading term. In Python, \suppfile{Python/numerical_checks.py}{mpmath} reproduced the grid and
\suppfile{Python/exact_arithmetic.py}{\texttt{fractions.Fraction}} checked the moment recurrences and rational
bounds. \suppfile{Python/identity_checks.py}{SymPy} \cite{SymPy2017} verified nine symbolic identities,
including the kernel decomposition and the Taylor coefficients of the phase.

\begin{table}[htbp]
\centering
\caption{\suppfile{Verification/table_1.csv}{Maxima of the scaled absolute error over the sampled $x$ values.}}
\label{tab:max}
\begin{tabular}{rrr}\hline
$q$ & $x$ at maximum & $\max q|R_q(x)|$ \\\hline
100 & -1 & 0.083681979528 \\
100.001 & -1 & 0.083681976829 \\
100.5 & -1 & 0.083680636233 \\
101 & -1 & 0.083679304576 \\
150 & -1 & 0.083587445357 \\
400 & -1 & 0.083445093725 \\
1600 & -1 & 0.083350517492 \\
\hline\end{tabular}
\end{table}

\begin{figure}[htbp]\centering
\setlength{\unitlength}{1pt}\begin{picture}(290,121)(-28,-24)
\put(0,0){\line(1,0){250}}\put(0,0){\line(0,1){85}}
\put(0,0){\line(0,-1){3}}\put(0,-11){\makebox(0,0){\scriptsize \(-1\)}}
\put(62.5,0){\line(0,-1){3}}\put(62.5,-11){\makebox(0,0){\scriptsize \(-0.5\)}}
\put(125,0){\line(0,-1){3}}\put(125,-11){\makebox(0,0){\scriptsize \(0\)}}
\put(187.5,0){\line(0,-1){3}}\put(187.5,-11){\makebox(0,0){\scriptsize \(0.5\)}}
\put(250,0){\line(0,-1){3}}\put(250,-11){\makebox(0,0){\scriptsize \(1\)}}
\put(0,0){\line(-1,0){3}}\put(-5,0){\makebox(0,0)[r]{\scriptsize \(0\)}}
\put(0,40){\line(-1,0){3}}\put(-5,40){\makebox(0,0)[r]{\scriptsize \(0.04\)}}
\put(0,80){\line(-1,0){3}}\put(-5,80){\makebox(0,0)[r]{\scriptsize \(0.08\)}}
\put(125,-21){\makebox(0,0){\small $x$}}\put(0,93){\makebox(0,0)[l]{\small $q|R_q(x)|$}}
\qbezier(0.000,83.6820)(3.125,81.7390)(6.250,79.7961)
\qbezier(6.250,79.7961)(9.375,77.8617)(12.500,75.9272)
\qbezier(12.500,75.9272)(15.625,74.0014)(18.750,72.0756)
\qbezier(18.750,72.0756)(21.875,70.1587)(25.000,68.2417)
\qbezier(25.000,68.2417)(28.125,66.3337)(31.250,64.4257)
\qbezier(31.250,64.4257)(34.375,62.5270)(37.500,60.6283)
\qbezier(37.500,60.6283)(40.625,58.7393)(43.750,56.8502)
\qbezier(43.750,56.8502)(46.875,54.9713)(50.000,53.0924)
\qbezier(50.000,53.0924)(53.125,51.2243)(56.250,49.3562)
\qbezier(56.250,49.3562)(59.375,47.5000)(62.500,45.6437)
\qbezier(62.500,45.6437)(65.625,43.8007)(68.750,41.9576)
\qbezier(68.750,41.9576)(71.875,40.1299)(75.000,38.3022)
\qbezier(75.000,38.3022)(78.125,36.4928)(81.250,34.6835)
\qbezier(81.250,34.6835)(84.375,32.8971)(87.500,31.1108)
\qbezier(87.500,31.1108)(90.625,29.3546)(93.750,27.5984)
\qbezier(93.750,27.5984)(96.875,25.8836)(100.000,24.1689)
\qbezier(100.000,24.1689)(103.125,22.5144)(106.250,20.8598)
\qbezier(106.250,20.8598)(109.375,19.2974)(112.500,17.7351)
\qbezier(112.500,17.7351)(115.625,16.3213)(118.750,14.9075)
\qbezier(118.750,14.9075)(121.875,13.7417)(125.000,12.5758)
\qbezier(125.000,12.5758)(128.125,11.8155)(131.250,11.0551)
\qbezier(131.250,11.0551)(134.375,10.8745)(137.500,10.6939)
\qbezier(137.500,10.6939)(140.625,11.1461)(143.750,11.5984)
\qbezier(143.750,11.5984)(146.875,12.5562)(150.000,13.5140)
\qbezier(150.000,13.5140)(153.125,14.7973)(156.250,16.0806)
\qbezier(156.250,16.0806)(159.375,17.5570)(162.500,19.0334)
\qbezier(162.500,19.0334)(165.625,20.6241)(168.750,22.2148)
\qbezier(168.750,22.2148)(171.875,23.8748)(175.000,25.5349)
\qbezier(175.000,25.5349)(178.125,27.2375)(181.250,28.9402)
\qbezier(181.250,28.9402)(184.375,30.6687)(187.500,32.3972)
\qbezier(187.500,32.3972)(190.625,34.1401)(193.750,35.8830)
\qbezier(193.750,35.8830)(196.875,37.6319)(200.000,39.3808)
\qbezier(200.000,39.3808)(203.125,41.1289)(206.250,42.8771)
\qbezier(206.250,42.8771)(209.375,44.6185)(212.500,46.3599)
\qbezier(212.500,46.3599)(215.625,48.0893)(218.750,49.8186)
\qbezier(218.750,49.8186)(221.875,51.5305)(225.000,53.2424)
\qbezier(225.000,53.2424)(228.125,54.9314)(231.250,56.6205)
\qbezier(231.250,56.6205)(234.375,58.2810)(237.500,59.9416)
\qbezier(237.500,59.9416)(240.625,61.5678)(243.750,63.1940)
\qbezier(243.750,63.1940)(246.875,64.7795)(250.000,66.3650)
\end{picture}
\caption{\suppfile{Figures/figure_1.pdf}{Scaled error at $q=100$.} \suppfile{Figures/figure_1_data.csv}{Consecutive samples are joined by
straight lines.}}\label{fig:slice}
\end{figure}

Figure~\ref{fig:correction} compares the absolute errors with one and two
terms. The correction reduces the error at all 123 points shown. At
$q=100$, the one-term error is approximately 59 to 471 times the
two-term error, depending on $x$.

\begin{figure}[htbp]
\centering
\setlength{\unitlength}{1pt}
\begin{picture}(416,148)
\linethickness{0.35pt}
\put(85.000,138.000){\makebox(0,0)[]{\scriptsize $q=100$}}
\qbezier(31.000,18.000)(31.000,71.000)(31.000,124.000)
\qbezier(31.000,18.000)(85.000,18.000)(139.000,18.000)
\qbezier(31.000,18.000)(31.000,17.000)(31.000,16.000)
\put(31.000,9.000){\makebox(0,0)[]{\scriptsize $-1$}}
\qbezier(85.000,18.000)(85.000,17.000)(85.000,16.000)
\put(85.000,9.000){\makebox(0,0)[]{\scriptsize $0$}}
\qbezier(139.000,18.000)(139.000,17.000)(139.000,16.000)
\put(139.000,9.000){\makebox(0,0)[]{\scriptsize $1$}}
\qbezier(31.000,123.750)(30.000,123.750)(29.000,123.750)
\put(27.000,123.750){\makebox(0,0)[r]{\scriptsize $0.1$}}
\qbezier(31.000,100.250)(30.000,100.250)(29.000,100.250)
\put(27.000,100.250){\makebox(0,0)[r]{\scriptsize $0.01$}}
\qbezier(31.000,76.750)(30.000,76.750)(29.000,76.750)
\put(27.000,76.750){\makebox(0,0)[r]{\scriptsize $0.001$}}
\qbezier(31.000,53.250)(30.000,53.250)(29.000,53.250)
\put(27.000,53.250){\makebox(0,0)[r]{\scriptsize $0.0001$}}
\qbezier(31.000,29.750)(30.000,29.750)(29.000,29.750)
\put(27.000,29.750){\makebox(0,0)[r]{\scriptsize $0.00001$}}
\qbezier(31.000,116.537)(31.770,116.541)(32.539,116.545)
\qbezier(33.700,116.551)(34.469,116.556)(35.239,116.560)
\qbezier(36.400,116.566)(37.169,116.570)(37.939,116.575)
\qbezier(39.100,116.581)(39.870,116.585)(40.639,116.589)
\qbezier(41.800,116.596)(42.569,116.600)(43.339,116.604)
\qbezier(44.500,116.610)(45.270,116.614)(46.039,116.619)
\qbezier(47.200,116.625)(47.969,116.629)(48.739,116.633)
\qbezier(49.900,116.639)(50.669,116.643)(51.439,116.647)
\qbezier(52.600,116.652)(53.370,116.656)(54.139,116.660)
\qbezier(55.300,116.666)(56.069,116.669)(56.839,116.673)
\qbezier(58.000,116.678)(58.770,116.682)(59.539,116.685)
\qbezier(60.700,116.690)(61.469,116.694)(62.239,116.697)
\qbezier(63.400,116.702)(64.169,116.705)(64.939,116.708)
\qbezier(66.100,116.712)(66.869,116.715)(67.639,116.718)
\qbezier(68.800,116.722)(69.570,116.724)(70.339,116.726)
\qbezier(71.500,116.730)(72.269,116.732)(73.039,116.734)
\qbezier(74.200,116.737)(74.970,116.739)(75.739,116.741)
\qbezier(76.900,116.744)(77.669,116.745)(78.439,116.746)
\qbezier(79.600,116.748)(80.369,116.749)(81.139,116.750)
\qbezier(82.300,116.751)(83.070,116.752)(83.839,116.752)
\qbezier(85.000,116.753)(85.769,116.753)(86.539,116.753)
\qbezier(87.700,116.753)(88.470,116.752)(89.239,116.752)
\qbezier(90.400,116.751)(91.169,116.750)(91.939,116.749)
\qbezier(93.100,116.747)(93.869,116.746)(94.639,116.744)
\qbezier(95.800,116.741)(96.570,116.739)(97.339,116.737)
\qbezier(98.500,116.733)(99.269,116.730)(100.039,116.728)
\qbezier(101.200,116.723)(101.970,116.720)(102.739,116.716)
\qbezier(103.900,116.711)(104.669,116.706)(105.439,116.702)
\qbezier(106.600,116.696)(107.369,116.691)(108.139,116.686)
\qbezier(109.300,116.679)(110.070,116.673)(110.839,116.667)
\qbezier(112.000,116.659)(112.769,116.653)(113.539,116.646)
\qbezier(114.700,116.637)(115.470,116.630)(116.239,116.623)
\qbezier(117.400,116.612)(118.169,116.604)(118.939,116.596)
\qbezier(120.100,116.585)(120.869,116.576)(121.639,116.568)
\qbezier(122.800,116.555)(123.570,116.546)(124.339,116.537)
\qbezier(125.500,116.523)(126.269,116.514)(127.039,116.504)
\qbezier(128.200,116.490)(128.969,116.479)(129.739,116.469)
\qbezier(130.900,116.454)(131.669,116.443)(132.439,116.432)
\qbezier(133.600,116.416)(134.370,116.405)(135.139,116.394)
\qbezier(136.300,116.377)(137.070,116.366)(137.839,116.354)
\put(30.350,115.887){\framebox(1.3,1.3){}}
\put(33.050,115.901){\framebox(1.3,1.3){}}
\put(35.750,115.916){\framebox(1.3,1.3){}}
\put(38.450,115.931){\framebox(1.3,1.3){}}
\put(41.150,115.946){\framebox(1.3,1.3){}}
\put(43.850,115.960){\framebox(1.3,1.3){}}
\put(46.550,115.975){\framebox(1.3,1.3){}}
\put(49.250,115.989){\framebox(1.3,1.3){}}
\put(51.950,116.002){\framebox(1.3,1.3){}}
\put(54.650,116.016){\framebox(1.3,1.3){}}
\put(57.350,116.028){\framebox(1.3,1.3){}}
\put(60.050,116.040){\framebox(1.3,1.3){}}
\put(62.750,116.052){\framebox(1.3,1.3){}}
\put(65.450,116.062){\framebox(1.3,1.3){}}
\put(68.150,116.072){\framebox(1.3,1.3){}}
\put(70.850,116.080){\framebox(1.3,1.3){}}
\put(73.550,116.087){\framebox(1.3,1.3){}}
\put(76.250,116.094){\framebox(1.3,1.3){}}
\put(78.950,116.098){\framebox(1.3,1.3){}}
\put(81.650,116.101){\framebox(1.3,1.3){}}
\put(84.350,116.103){\framebox(1.3,1.3){}}
\put(87.050,116.103){\framebox(1.3,1.3){}}
\put(89.750,116.101){\framebox(1.3,1.3){}}
\put(92.450,116.097){\framebox(1.3,1.3){}}
\put(95.150,116.091){\framebox(1.3,1.3){}}
\put(97.850,116.083){\framebox(1.3,1.3){}}
\put(100.550,116.073){\framebox(1.3,1.3){}}
\put(103.250,116.061){\framebox(1.3,1.3){}}
\put(105.950,116.046){\framebox(1.3,1.3){}}
\put(108.650,116.029){\framebox(1.3,1.3){}}
\put(111.350,116.009){\framebox(1.3,1.3){}}
\put(114.050,115.987){\framebox(1.3,1.3){}}
\put(116.750,115.962){\framebox(1.3,1.3){}}
\put(119.450,115.935){\framebox(1.3,1.3){}}
\put(122.150,115.905){\framebox(1.3,1.3){}}
\put(124.850,115.873){\framebox(1.3,1.3){}}
\put(127.550,115.840){\framebox(1.3,1.3){}}
\put(130.250,115.804){\framebox(1.3,1.3){}}
\put(132.950,115.766){\framebox(1.3,1.3){}}
\put(135.650,115.727){\framebox(1.3,1.3){}}
\put(138.350,115.687){\framebox(1.3,1.3){}}
\qbezier(31.000,74.932)(32.350,74.689)(33.700,74.447)
\qbezier(33.700,74.447)(35.050,74.193)(36.400,73.939)
\qbezier(36.400,73.939)(37.750,73.674)(39.100,73.408)
\qbezier(39.100,73.408)(40.450,73.129)(41.800,72.850)
\qbezier(41.800,72.850)(43.150,72.557)(44.500,72.263)
\qbezier(44.500,72.263)(45.850,71.953)(47.200,71.643)
\qbezier(47.200,71.643)(48.550,71.315)(49.900,70.986)
\qbezier(49.900,70.986)(51.250,70.637)(52.600,70.288)
\qbezier(52.600,70.288)(53.950,69.916)(55.300,69.544)
\qbezier(55.300,69.544)(56.650,69.144)(58.000,68.745)
\qbezier(58.000,68.745)(59.350,68.316)(60.700,67.886)
\qbezier(60.700,67.886)(62.050,67.421)(63.400,66.956)
\qbezier(63.400,66.956)(64.750,66.449)(66.100,65.943)
\qbezier(66.100,65.943)(67.450,65.388)(68.800,64.833)
\qbezier(68.800,64.833)(70.150,64.222)(71.500,63.611)
\qbezier(71.500,63.611)(72.850,62.934)(74.200,62.257)
\qbezier(74.200,62.257)(75.550,61.505)(76.900,60.754)
\qbezier(76.900,60.754)(78.250,59.926)(79.600,59.098)
\qbezier(79.600,59.098)(80.950,58.211)(82.300,57.325)
\qbezier(82.300,57.325)(83.650,56.457)(85.000,55.589)
\qbezier(85.000,55.589)(86.350,54.931)(87.700,54.274)
\qbezier(87.700,54.274)(89.050,54.104)(90.400,53.935)
\qbezier(90.400,53.935)(91.750,54.349)(93.100,54.763)
\qbezier(93.100,54.763)(94.450,55.543)(95.800,56.323)
\qbezier(95.800,56.323)(97.150,57.211)(98.500,58.098)
\qbezier(98.500,58.098)(99.850,58.958)(101.200,59.819)
\qbezier(101.200,59.819)(102.550,60.607)(103.900,61.396)
\qbezier(103.900,61.396)(105.250,62.107)(106.600,62.818)
\qbezier(106.600,62.818)(107.950,63.456)(109.300,64.095)
\qbezier(109.300,64.095)(110.650,64.671)(112.000,65.247)
\qbezier(112.000,65.247)(113.350,65.768)(114.700,66.290)
\qbezier(114.700,66.290)(116.050,66.765)(117.400,67.239)
\qbezier(117.400,67.239)(118.750,67.673)(120.100,68.107)
\qbezier(120.100,68.107)(121.450,68.506)(122.800,68.904)
\qbezier(122.800,68.904)(124.150,69.272)(125.500,69.639)
\qbezier(125.500,69.639)(126.850,69.978)(128.200,70.317)
\qbezier(128.200,70.317)(129.550,70.631)(130.900,70.945)
\qbezier(130.900,70.945)(132.250,71.236)(133.600,71.527)
\qbezier(133.600,71.527)(134.950,71.796)(136.300,72.066)
\qbezier(136.300,72.066)(137.650,72.316)(139.000,72.566)
\put(30.350,74.282){\rule{1.3pt}{1.3pt}}
\put(33.050,73.797){\rule{1.3pt}{1.3pt}}
\put(35.750,73.289){\rule{1.3pt}{1.3pt}}
\put(38.450,72.758){\rule{1.3pt}{1.3pt}}
\put(41.150,72.200){\rule{1.3pt}{1.3pt}}
\put(43.850,71.613){\rule{1.3pt}{1.3pt}}
\put(46.550,70.993){\rule{1.3pt}{1.3pt}}
\put(49.250,70.336){\rule{1.3pt}{1.3pt}}
\put(51.950,69.638){\rule{1.3pt}{1.3pt}}
\put(54.650,68.894){\rule{1.3pt}{1.3pt}}
\put(57.350,68.095){\rule{1.3pt}{1.3pt}}
\put(60.050,67.236){\rule{1.3pt}{1.3pt}}
\put(62.750,66.306){\rule{1.3pt}{1.3pt}}
\put(65.450,65.293){\rule{1.3pt}{1.3pt}}
\put(68.150,64.183){\rule{1.3pt}{1.3pt}}
\put(70.850,62.961){\rule{1.3pt}{1.3pt}}
\put(73.550,61.607){\rule{1.3pt}{1.3pt}}
\put(76.250,60.104){\rule{1.3pt}{1.3pt}}
\put(78.950,58.448){\rule{1.3pt}{1.3pt}}
\put(81.650,56.675){\rule{1.3pt}{1.3pt}}
\put(84.350,54.939){\rule{1.3pt}{1.3pt}}
\put(87.050,53.624){\rule{1.3pt}{1.3pt}}
\put(89.750,53.285){\rule{1.3pt}{1.3pt}}
\put(92.450,54.113){\rule{1.3pt}{1.3pt}}
\put(95.150,55.673){\rule{1.3pt}{1.3pt}}
\put(97.850,57.448){\rule{1.3pt}{1.3pt}}
\put(100.550,59.169){\rule{1.3pt}{1.3pt}}
\put(103.250,60.746){\rule{1.3pt}{1.3pt}}
\put(105.950,62.168){\rule{1.3pt}{1.3pt}}
\put(108.650,63.445){\rule{1.3pt}{1.3pt}}
\put(111.350,64.597){\rule{1.3pt}{1.3pt}}
\put(114.050,65.640){\rule{1.3pt}{1.3pt}}
\put(116.750,66.589){\rule{1.3pt}{1.3pt}}
\put(119.450,67.457){\rule{1.3pt}{1.3pt}}
\put(122.150,68.254){\rule{1.3pt}{1.3pt}}
\put(124.850,68.989){\rule{1.3pt}{1.3pt}}
\put(127.550,69.667){\rule{1.3pt}{1.3pt}}
\put(130.250,70.295){\rule{1.3pt}{1.3pt}}
\put(132.950,70.877){\rule{1.3pt}{1.3pt}}
\put(135.650,71.416){\rule{1.3pt}{1.3pt}}
\put(138.350,71.916){\rule{1.3pt}{1.3pt}}
\put(222.000,138.000){\makebox(0,0)[]{\scriptsize $q=400$}}
\qbezier(168.000,18.000)(168.000,71.000)(168.000,124.000)
\qbezier(168.000,18.000)(222.000,18.000)(276.000,18.000)
\qbezier(168.000,18.000)(168.000,17.000)(168.000,16.000)
\put(168.000,9.000){\makebox(0,0)[]{\scriptsize $-1$}}
\qbezier(222.000,18.000)(222.000,17.000)(222.000,16.000)
\put(222.000,9.000){\makebox(0,0)[]{\scriptsize $0$}}
\qbezier(276.000,18.000)(276.000,17.000)(276.000,16.000)
\put(276.000,9.000){\makebox(0,0)[]{\scriptsize $1$}}
\qbezier(168.000,123.750)(167.000,123.750)(166.000,123.750)
\qbezier(168.000,100.250)(167.000,100.250)(166.000,100.250)
\qbezier(168.000,76.750)(167.000,76.750)(166.000,76.750)
\qbezier(168.000,53.250)(167.000,53.250)(166.000,53.250)
\qbezier(168.000,29.750)(167.000,29.750)(166.000,29.750)
\qbezier(168.000,109.459)(168.769,109.463)(169.539,109.467)
\qbezier(170.700,109.473)(171.469,109.477)(172.239,109.481)
\qbezier(173.400,109.487)(174.169,109.491)(174.939,109.496)
\qbezier(176.100,109.502)(176.870,109.506)(177.639,109.510)
\qbezier(178.800,109.516)(179.570,109.520)(180.339,109.524)
\qbezier(181.500,109.530)(182.269,109.534)(183.039,109.538)
\qbezier(184.200,109.544)(184.969,109.548)(185.739,109.552)
\qbezier(186.900,109.558)(187.669,109.562)(188.439,109.565)
\qbezier(189.600,109.571)(190.370,109.575)(191.139,109.578)
\qbezier(192.300,109.584)(193.070,109.587)(193.839,109.591)
\qbezier(195.000,109.596)(195.769,109.599)(196.539,109.603)
\qbezier(197.700,109.608)(198.469,109.611)(199.239,109.614)
\qbezier(200.400,109.619)(201.169,109.621)(201.939,109.624)
\qbezier(203.100,109.629)(203.870,109.631)(204.639,109.634)
\qbezier(205.800,109.638)(206.570,109.640)(207.339,109.642)
\qbezier(208.500,109.646)(209.269,109.648)(210.039,109.650)
\qbezier(211.200,109.653)(211.969,109.654)(212.739,109.656)
\qbezier(213.900,109.658)(214.669,109.660)(215.439,109.661)
\qbezier(216.600,109.663)(217.370,109.663)(218.139,109.664)
\qbezier(219.300,109.666)(220.070,109.666)(220.839,109.666)
\qbezier(222.000,109.667)(222.769,109.667)(223.539,109.666)
\qbezier(224.700,109.666)(225.469,109.666)(226.239,109.665)
\qbezier(227.400,109.664)(228.169,109.663)(228.939,109.662)
\qbezier(230.100,109.660)(230.870,109.658)(231.639,109.656)
\qbezier(232.800,109.653)(233.570,109.651)(234.339,109.649)
\qbezier(235.500,109.645)(236.269,109.642)(237.039,109.639)
\qbezier(238.200,109.634)(238.969,109.631)(239.739,109.627)
\qbezier(240.900,109.622)(241.669,109.617)(242.439,109.613)
\qbezier(243.600,109.606)(244.370,109.601)(245.139,109.596)
\qbezier(246.300,109.588)(247.070,109.583)(247.839,109.577)
\qbezier(249.000,109.568)(249.769,109.562)(250.539,109.555)
\qbezier(251.700,109.545)(252.469,109.538)(253.239,109.531)
\qbezier(254.400,109.520)(255.170,109.512)(255.939,109.504)
\qbezier(257.100,109.492)(257.870,109.484)(258.639,109.475)
\qbezier(259.800,109.462)(260.570,109.453)(261.339,109.443)
\qbezier(262.500,109.429)(263.269,109.419)(264.039,109.409)
\qbezier(265.200,109.394)(265.969,109.384)(266.739,109.373)
\qbezier(267.900,109.358)(268.669,109.347)(269.439,109.336)
\qbezier(270.600,109.319)(271.370,109.308)(272.139,109.296)
\qbezier(273.300,109.279)(274.070,109.267)(274.839,109.255)
\put(167.350,108.809){\framebox(1.3,1.3){}}
\put(170.050,108.823){\framebox(1.3,1.3){}}
\put(172.750,108.837){\framebox(1.3,1.3){}}
\put(175.450,108.852){\framebox(1.3,1.3){}}
\put(178.150,108.866){\framebox(1.3,1.3){}}
\put(180.850,108.880){\framebox(1.3,1.3){}}
\put(183.550,108.894){\framebox(1.3,1.3){}}
\put(186.250,108.908){\framebox(1.3,1.3){}}
\put(188.950,108.921){\framebox(1.3,1.3){}}
\put(191.650,108.934){\framebox(1.3,1.3){}}
\put(194.350,108.946){\framebox(1.3,1.3){}}
\put(197.050,108.958){\framebox(1.3,1.3){}}
\put(199.750,108.969){\framebox(1.3,1.3){}}
\put(202.450,108.979){\framebox(1.3,1.3){}}
\put(205.150,108.988){\framebox(1.3,1.3){}}
\put(207.850,108.996){\framebox(1.3,1.3){}}
\put(210.550,109.003){\framebox(1.3,1.3){}}
\put(213.250,109.008){\framebox(1.3,1.3){}}
\put(215.950,109.013){\framebox(1.3,1.3){}}
\put(218.650,109.016){\framebox(1.3,1.3){}}
\put(221.350,109.017){\framebox(1.3,1.3){}}
\put(224.050,109.016){\framebox(1.3,1.3){}}
\put(226.750,109.014){\framebox(1.3,1.3){}}
\put(229.450,109.010){\framebox(1.3,1.3){}}
\put(232.150,109.003){\framebox(1.3,1.3){}}
\put(234.850,108.995){\framebox(1.3,1.3){}}
\put(237.550,108.984){\framebox(1.3,1.3){}}
\put(240.250,108.972){\framebox(1.3,1.3){}}
\put(242.950,108.956){\framebox(1.3,1.3){}}
\put(245.650,108.938){\framebox(1.3,1.3){}}
\put(248.350,108.918){\framebox(1.3,1.3){}}
\put(251.050,108.895){\framebox(1.3,1.3){}}
\put(253.750,108.870){\framebox(1.3,1.3){}}
\put(256.450,108.842){\framebox(1.3,1.3){}}
\put(259.150,108.812){\framebox(1.3,1.3){}}
\put(261.850,108.779){\framebox(1.3,1.3){}}
\put(264.550,108.744){\framebox(1.3,1.3){}}
\put(267.250,108.708){\framebox(1.3,1.3){}}
\put(269.950,108.669){\framebox(1.3,1.3){}}
\put(272.650,108.629){\framebox(1.3,1.3){}}
\put(275.350,108.588){\framebox(1.3,1.3){}}
\qbezier(168.000,60.755)(169.350,60.510)(170.700,60.266)
\qbezier(170.700,60.266)(172.050,60.011)(173.400,59.755)
\qbezier(173.400,59.755)(174.750,59.487)(176.100,59.220)
\qbezier(176.100,59.220)(177.450,58.938)(178.800,58.657)
\qbezier(178.800,58.657)(180.150,58.360)(181.500,58.064)
\qbezier(181.500,58.064)(182.850,57.750)(184.200,57.437)
\qbezier(184.200,57.437)(185.550,57.104)(186.900,56.772)
\qbezier(186.900,56.772)(188.250,56.418)(189.600,56.064)
\qbezier(189.600,56.064)(190.950,55.686)(192.300,55.308)
\qbezier(192.300,55.308)(193.650,54.901)(195.000,54.495)
\qbezier(195.000,54.495)(196.350,54.056)(197.700,53.616)
\qbezier(197.700,53.616)(199.050,53.139)(200.400,52.662)
\qbezier(200.400,52.662)(201.750,52.140)(203.100,51.618)
\qbezier(203.100,51.618)(204.450,51.041)(205.800,50.465)
\qbezier(205.800,50.465)(207.150,49.824)(208.500,49.183)
\qbezier(208.500,49.183)(209.850,48.464)(211.200,47.744)
\qbezier(211.200,47.744)(212.550,46.929)(213.900,46.113)
\qbezier(213.900,46.113)(215.250,45.187)(216.600,44.260)
\qbezier(216.600,44.260)(217.950,43.219)(219.300,42.179)
\qbezier(219.300,42.179)(220.650,41.082)(222.000,39.985)
\qbezier(222.000,39.985)(223.350,39.071)(224.700,38.158)
\qbezier(224.700,38.158)(226.050,37.907)(227.400,37.657)
\qbezier(227.400,37.657)(228.750,38.262)(230.100,38.867)
\qbezier(230.100,38.867)(231.450,39.906)(232.800,40.945)
\qbezier(232.800,40.945)(234.150,42.026)(235.500,43.107)
\qbezier(235.500,43.107)(236.850,44.092)(238.200,45.077)
\qbezier(238.200,45.077)(239.550,45.945)(240.900,46.813)
\qbezier(240.900,46.813)(242.250,47.575)(243.600,48.337)
\qbezier(243.600,48.337)(244.950,49.010)(246.300,49.684)
\qbezier(246.300,49.684)(247.650,50.284)(249.000,50.883)
\qbezier(249.000,50.883)(250.350,51.422)(251.700,51.960)
\qbezier(251.700,51.960)(253.050,52.447)(254.400,52.935)
\qbezier(254.400,52.935)(255.750,53.378)(257.100,53.821)
\qbezier(257.100,53.821)(258.450,54.226)(259.800,54.632)
\qbezier(259.800,54.632)(261.150,55.004)(262.500,55.377)
\qbezier(262.500,55.377)(263.850,55.720)(265.200,56.063)
\qbezier(265.200,56.063)(266.550,56.380)(267.900,56.697)
\qbezier(267.900,56.697)(269.250,56.990)(270.600,57.283)
\qbezier(270.600,57.283)(271.950,57.554)(273.300,57.825)
\qbezier(273.300,57.825)(274.650,58.076)(276.000,58.327)
\put(167.350,60.105){\rule{1.3pt}{1.3pt}}
\put(170.050,59.616){\rule{1.3pt}{1.3pt}}
\put(172.750,59.105){\rule{1.3pt}{1.3pt}}
\put(175.450,58.570){\rule{1.3pt}{1.3pt}}
\put(178.150,58.007){\rule{1.3pt}{1.3pt}}
\put(180.850,57.414){\rule{1.3pt}{1.3pt}}
\put(183.550,56.787){\rule{1.3pt}{1.3pt}}
\put(186.250,56.122){\rule{1.3pt}{1.3pt}}
\put(188.950,55.414){\rule{1.3pt}{1.3pt}}
\put(191.650,54.658){\rule{1.3pt}{1.3pt}}
\put(194.350,53.845){\rule{1.3pt}{1.3pt}}
\put(197.050,52.966){\rule{1.3pt}{1.3pt}}
\put(199.750,52.012){\rule{1.3pt}{1.3pt}}
\put(202.450,50.968){\rule{1.3pt}{1.3pt}}
\put(205.150,49.815){\rule{1.3pt}{1.3pt}}
\put(207.850,48.533){\rule{1.3pt}{1.3pt}}
\put(210.550,47.094){\rule{1.3pt}{1.3pt}}
\put(213.250,45.463){\rule{1.3pt}{1.3pt}}
\put(215.950,43.610){\rule{1.3pt}{1.3pt}}
\put(218.650,41.529){\rule{1.3pt}{1.3pt}}
\put(221.350,39.335){\rule{1.3pt}{1.3pt}}
\put(224.050,37.508){\rule{1.3pt}{1.3pt}}
\put(226.750,37.007){\rule{1.3pt}{1.3pt}}
\put(229.450,38.217){\rule{1.3pt}{1.3pt}}
\put(232.150,40.295){\rule{1.3pt}{1.3pt}}
\put(234.850,42.457){\rule{1.3pt}{1.3pt}}
\put(237.550,44.427){\rule{1.3pt}{1.3pt}}
\put(240.250,46.163){\rule{1.3pt}{1.3pt}}
\put(242.950,47.687){\rule{1.3pt}{1.3pt}}
\put(245.650,49.034){\rule{1.3pt}{1.3pt}}
\put(248.350,50.233){\rule{1.3pt}{1.3pt}}
\put(251.050,51.310){\rule{1.3pt}{1.3pt}}
\put(253.750,52.285){\rule{1.3pt}{1.3pt}}
\put(256.450,53.171){\rule{1.3pt}{1.3pt}}
\put(259.150,53.982){\rule{1.3pt}{1.3pt}}
\put(261.850,54.727){\rule{1.3pt}{1.3pt}}
\put(264.550,55.413){\rule{1.3pt}{1.3pt}}
\put(267.250,56.047){\rule{1.3pt}{1.3pt}}
\put(269.950,56.633){\rule{1.3pt}{1.3pt}}
\put(272.650,57.175){\rule{1.3pt}{1.3pt}}
\put(275.350,57.677){\rule{1.3pt}{1.3pt}}
\put(359.000,138.000){\makebox(0,0)[]{\scriptsize $q=1600$}}
\qbezier(305.000,18.000)(305.000,71.000)(305.000,124.000)
\qbezier(305.000,18.000)(359.000,18.000)(413.000,18.000)
\qbezier(305.000,18.000)(305.000,17.000)(305.000,16.000)
\put(305.000,9.000){\makebox(0,0)[]{\scriptsize $-1$}}
\qbezier(359.000,18.000)(359.000,17.000)(359.000,16.000)
\put(359.000,9.000){\makebox(0,0)[]{\scriptsize $0$}}
\qbezier(413.000,18.000)(413.000,17.000)(413.000,16.000)
\put(413.000,9.000){\makebox(0,0)[]{\scriptsize $1$}}
\qbezier(305.000,123.750)(304.000,123.750)(303.000,123.750)
\qbezier(305.000,100.250)(304.000,100.250)(303.000,100.250)
\qbezier(305.000,76.750)(304.000,76.750)(303.000,76.750)
\qbezier(305.000,53.250)(304.000,53.250)(303.000,53.250)
\qbezier(305.000,29.750)(304.000,29.750)(303.000,29.750)
\qbezier(305.000,102.384)(305.769,102.388)(306.539,102.392)
\qbezier(307.700,102.398)(308.469,102.402)(309.239,102.407)
\qbezier(310.400,102.413)(311.169,102.417)(311.939,102.421)
\qbezier(313.100,102.427)(313.870,102.431)(314.639,102.435)
\qbezier(315.800,102.441)(316.570,102.445)(317.339,102.449)
\qbezier(318.500,102.455)(319.269,102.459)(320.039,102.463)
\qbezier(321.200,102.468)(321.969,102.472)(322.739,102.476)
\qbezier(323.900,102.482)(324.669,102.486)(325.439,102.489)
\qbezier(326.600,102.495)(327.370,102.498)(328.139,102.502)
\qbezier(329.300,102.507)(330.070,102.511)(330.839,102.514)
\qbezier(332.000,102.519)(332.769,102.523)(333.539,102.526)
\qbezier(334.700,102.531)(335.469,102.534)(336.239,102.537)
\qbezier(337.400,102.541)(338.169,102.544)(338.939,102.547)
\qbezier(340.100,102.551)(340.870,102.554)(341.639,102.556)
\qbezier(342.800,102.560)(343.570,102.562)(344.339,102.565)
\qbezier(345.500,102.568)(346.269,102.570)(347.039,102.572)
\qbezier(348.200,102.575)(348.969,102.576)(349.739,102.578)
\qbezier(350.900,102.580)(351.669,102.581)(352.439,102.583)
\qbezier(353.600,102.584)(354.370,102.585)(355.139,102.586)
\qbezier(356.300,102.587)(357.070,102.587)(357.839,102.588)
\qbezier(359.000,102.588)(359.769,102.588)(360.539,102.588)
\qbezier(361.700,102.587)(362.469,102.587)(363.239,102.586)
\qbezier(364.400,102.585)(365.169,102.584)(365.939,102.582)
\qbezier(367.100,102.580)(367.870,102.579)(368.639,102.577)
\qbezier(369.800,102.574)(370.570,102.572)(371.339,102.569)
\qbezier(372.500,102.565)(373.269,102.562)(374.039,102.559)
\qbezier(375.200,102.555)(375.969,102.551)(376.739,102.547)
\qbezier(377.900,102.541)(378.669,102.537)(379.439,102.533)
\qbezier(380.600,102.526)(381.370,102.521)(382.139,102.516)
\qbezier(383.300,102.508)(384.070,102.502)(384.839,102.496)
\qbezier(386.000,102.487)(386.769,102.481)(387.539,102.474)
\qbezier(388.700,102.464)(389.469,102.457)(390.239,102.450)
\qbezier(391.400,102.439)(392.169,102.431)(392.939,102.423)
\qbezier(394.100,102.410)(394.870,102.402)(395.639,102.393)
\qbezier(396.800,102.380)(397.570,102.370)(398.339,102.361)
\qbezier(399.500,102.347)(400.269,102.337)(401.039,102.327)
\qbezier(402.200,102.312)(402.969,102.301)(403.739,102.290)
\qbezier(404.900,102.274)(405.669,102.263)(406.439,102.252)
\qbezier(407.600,102.235)(408.370,102.224)(409.139,102.212)
\qbezier(410.300,102.194)(411.070,102.183)(411.839,102.171)
\put(304.350,101.734){\framebox(1.3,1.3){}}
\put(307.050,101.748){\framebox(1.3,1.3){}}
\put(309.750,101.763){\framebox(1.3,1.3){}}
\put(312.450,101.777){\framebox(1.3,1.3){}}
\put(315.150,101.791){\framebox(1.3,1.3){}}
\put(317.850,101.805){\framebox(1.3,1.3){}}
\put(320.550,101.818){\framebox(1.3,1.3){}}
\put(323.250,101.832){\framebox(1.3,1.3){}}
\put(325.950,101.845){\framebox(1.3,1.3){}}
\put(328.650,101.857){\framebox(1.3,1.3){}}
\put(331.350,101.869){\framebox(1.3,1.3){}}
\put(334.050,101.881){\framebox(1.3,1.3){}}
\put(336.750,101.891){\framebox(1.3,1.3){}}
\put(339.450,101.901){\framebox(1.3,1.3){}}
\put(342.150,101.910){\framebox(1.3,1.3){}}
\put(344.850,101.918){\framebox(1.3,1.3){}}
\put(347.550,101.925){\framebox(1.3,1.3){}}
\put(350.250,101.930){\framebox(1.3,1.3){}}
\put(352.950,101.934){\framebox(1.3,1.3){}}
\put(355.650,101.937){\framebox(1.3,1.3){}}
\put(358.350,101.938){\framebox(1.3,1.3){}}
\put(361.050,101.937){\framebox(1.3,1.3){}}
\put(363.750,101.935){\framebox(1.3,1.3){}}
\put(366.450,101.930){\framebox(1.3,1.3){}}
\put(369.150,101.924){\framebox(1.3,1.3){}}
\put(371.850,101.915){\framebox(1.3,1.3){}}
\put(374.550,101.905){\framebox(1.3,1.3){}}
\put(377.250,101.891){\framebox(1.3,1.3){}}
\put(379.950,101.876){\framebox(1.3,1.3){}}
\put(382.650,101.858){\framebox(1.3,1.3){}}
\put(385.350,101.837){\framebox(1.3,1.3){}}
\put(388.050,101.814){\framebox(1.3,1.3){}}
\put(390.750,101.789){\framebox(1.3,1.3){}}
\put(393.450,101.760){\framebox(1.3,1.3){}}
\put(396.150,101.730){\framebox(1.3,1.3){}}
\put(398.850,101.697){\framebox(1.3,1.3){}}
\put(401.550,101.662){\framebox(1.3,1.3){}}
\put(404.250,101.624){\framebox(1.3,1.3){}}
\put(406.950,101.585){\framebox(1.3,1.3){}}
\put(409.650,101.544){\framebox(1.3,1.3){}}
\put(412.350,101.503){\framebox(1.3,1.3){}}
\qbezier(305.000,46.595)(306.350,46.350)(307.700,46.105)
\qbezier(307.700,46.105)(309.050,45.848)(310.400,45.592)
\qbezier(310.400,45.592)(311.750,45.323)(313.100,45.055)
\qbezier(313.100,45.055)(314.450,44.772)(315.800,44.490)
\qbezier(315.800,44.490)(317.150,44.192)(318.500,43.894)
\qbezier(318.500,43.894)(319.850,43.579)(321.200,43.264)
\qbezier(321.200,43.264)(322.550,42.930)(323.900,42.596)
\qbezier(323.900,42.596)(325.250,42.240)(326.600,41.884)
\qbezier(326.600,41.884)(327.950,41.503)(329.300,41.122)
\qbezier(329.300,41.122)(330.650,40.712)(332.000,40.303)
\qbezier(332.000,40.303)(333.350,39.860)(334.700,39.416)
\qbezier(334.700,39.416)(336.050,38.934)(337.400,38.451)
\qbezier(337.400,38.451)(338.750,37.922)(340.100,37.393)
\qbezier(340.100,37.393)(341.450,36.807)(342.800,36.222)
\qbezier(342.800,36.222)(344.150,35.567)(345.500,34.913)
\qbezier(345.500,34.913)(346.850,34.174)(348.200,33.434)
\qbezier(348.200,33.434)(349.550,32.589)(350.900,31.743)
\qbezier(350.900,31.743)(352.250,30.768)(353.600,29.792)
\qbezier(353.600,29.792)(354.950,28.668)(356.300,27.545)
\qbezier(356.300,27.545)(357.650,26.309)(359.000,25.074)
\qbezier(359.000,25.074)(360.350,23.979)(361.700,22.885)
\qbezier(361.700,22.885)(363.050,22.570)(364.400,22.255)
\qbezier(364.400,22.255)(365.750,23.005)(367.100,23.756)
\qbezier(367.100,23.756)(368.450,24.965)(369.800,26.173)
\qbezier(369.800,26.173)(371.150,27.365)(372.500,28.556)
\qbezier(372.500,28.556)(373.850,29.606)(375.200,30.655)
\qbezier(375.200,30.655)(376.550,31.561)(377.900,32.467)
\qbezier(377.900,32.467)(379.250,33.254)(380.600,34.040)
\qbezier(380.600,34.040)(381.950,34.729)(383.300,35.418)
\qbezier(383.300,35.418)(384.650,36.029)(386.000,36.639)
\qbezier(386.000,36.639)(387.350,37.185)(388.700,37.731)
\qbezier(388.700,37.731)(390.050,38.224)(391.400,38.716)
\qbezier(391.400,38.716)(392.750,39.163)(394.100,39.611)
\qbezier(394.100,39.611)(395.450,40.019)(396.800,40.428)
\qbezier(396.800,40.428)(398.150,40.802)(399.500,41.177)
\qbezier(399.500,41.177)(400.850,41.521)(402.200,41.866)
\qbezier(402.200,41.866)(403.550,42.184)(404.900,42.502)
\qbezier(404.900,42.502)(406.250,42.796)(407.600,43.090)
\qbezier(407.600,43.090)(408.950,43.362)(410.300,43.634)
\qbezier(410.300,43.634)(411.650,43.885)(413.000,44.136)
\put(304.350,45.945){\rule{1.3pt}{1.3pt}}
\put(307.050,45.455){\rule{1.3pt}{1.3pt}}
\put(309.750,44.942){\rule{1.3pt}{1.3pt}}
\put(312.450,44.405){\rule{1.3pt}{1.3pt}}
\put(315.150,43.840){\rule{1.3pt}{1.3pt}}
\put(317.850,43.244){\rule{1.3pt}{1.3pt}}
\put(320.550,42.614){\rule{1.3pt}{1.3pt}}
\put(323.250,41.946){\rule{1.3pt}{1.3pt}}
\put(325.950,41.234){\rule{1.3pt}{1.3pt}}
\put(328.650,40.472){\rule{1.3pt}{1.3pt}}
\put(331.350,39.653){\rule{1.3pt}{1.3pt}}
\put(334.050,38.766){\rule{1.3pt}{1.3pt}}
\put(336.750,37.801){\rule{1.3pt}{1.3pt}}
\put(339.450,36.743){\rule{1.3pt}{1.3pt}}
\put(342.150,35.572){\rule{1.3pt}{1.3pt}}
\put(344.850,34.263){\rule{1.3pt}{1.3pt}}
\put(347.550,32.784){\rule{1.3pt}{1.3pt}}
\put(350.250,31.093){\rule{1.3pt}{1.3pt}}
\put(352.950,29.142){\rule{1.3pt}{1.3pt}}
\put(355.650,26.895){\rule{1.3pt}{1.3pt}}
\put(358.350,24.424){\rule{1.3pt}{1.3pt}}
\put(361.050,22.235){\rule{1.3pt}{1.3pt}}
\put(363.750,21.605){\rule{1.3pt}{1.3pt}}
\put(366.450,23.106){\rule{1.3pt}{1.3pt}}
\put(369.150,25.523){\rule{1.3pt}{1.3pt}}
\put(371.850,27.906){\rule{1.3pt}{1.3pt}}
\put(374.550,30.005){\rule{1.3pt}{1.3pt}}
\put(377.250,31.817){\rule{1.3pt}{1.3pt}}
\put(379.950,33.390){\rule{1.3pt}{1.3pt}}
\put(382.650,34.768){\rule{1.3pt}{1.3pt}}
\put(385.350,35.989){\rule{1.3pt}{1.3pt}}
\put(388.050,37.081){\rule{1.3pt}{1.3pt}}
\put(390.750,38.066){\rule{1.3pt}{1.3pt}}
\put(393.450,38.961){\rule{1.3pt}{1.3pt}}
\put(396.150,39.778){\rule{1.3pt}{1.3pt}}
\put(398.850,40.527){\rule{1.3pt}{1.3pt}}
\put(401.550,41.216){\rule{1.3pt}{1.3pt}}
\put(404.250,41.852){\rule{1.3pt}{1.3pt}}
\put(406.950,42.440){\rule{1.3pt}{1.3pt}}
\put(409.650,42.984){\rule{1.3pt}{1.3pt}}
\put(412.350,43.486){\rule{1.3pt}{1.3pt}}
\end{picture}
\caption{\suppfile{Figures/figure_2.pdf}{Absolute error as a function of $x$ at the} \suppfile{Figures/figure_2_data.csv}{41 grid points} in
each panel. The dashed curve with open markers shows $|Z_q(x)-F(x)|$.
The solid curve with filled markers shows $|R_q(x)|$ after adding
$G(x)/\sqrt{q}$. All panels use the same logarithmic vertical scale.}
\label{fig:correction}
\end{figure}

In Mathematica we also integrated the normalized form of \eqref{eq:ap}
over intervals with endpoints $0,1,\sqrt q,4\sqrt q,\infty$, using 60 working
digits and accuracy and precision goals of 35 digits. The quotient at zero
has the limiting value $-\ii s/(\pi q)$. For $q=100$ and $x=-1,0,1$, the
absolute differences from the 70-digit Hurwitz zeta evaluation were about
$7.53\times10^{-58}$, $1.67\times10^{-54}$ and $1.19\times10^{-53}$,
respectively. The integral bounds give the uniform estimate
in Theorem~\ref{thm:main}.

\section{Conclusion}
After the correction of order $q^{-1/2}$, the Taylor remainder and the
residual integral together give the proposed bound $4.3/q$, including
at the endpoints. The whole-axis estimate controls the residual integral
without a splitting point.

The optimal constant remains undetermined. Comparing constants with Lerch
integral formulas and approximate functional equations requires a derivation
in the same normalization and parameter range.
\appendix
\section*{Appendix 1}
\addcontentsline{toc}{section}{Appendix 1}
\enlargethispage{\baselineskip}

Earlier working bounds were $24.8/q$ and $11.3/q$. For the short-interval
multiplier $3$, the calculation in Lemma~\ref{lem:local} yields
\[
 B_3=\frac3{12}+\frac3{240}+\frac{22}{7\cdot120000}+10^{-5},
 \qquad
 \overline J/2+B_3=\frac{356452924783991052355217}{86597871701497153200000}<4.117 .
\]
We retain this as a possible refinement pending separate validation. Equation~\eqref{eq:otherC} comes from a different kernel
majorant and gives a slightly different rational constant.

The short-interval multiplier cannot be used on the whole ray. At
$q=y=100$, $x=0$, set $A=100\pi(1-\pi/4)>64$. Then
$|w|=2^{-1/4}\e^{-A}$ and $D=2A+\ii\pi/4$. Consequently
$|w(1-\e^D)|\ge 2^{-1/4}(\e^A-\e^{-A})>5|D|$,
using $\e^A>A^3/6$ and $5|D|<3150$. Formula~\eqref{eq:whole}
continues to hold at this point.

A second way to check the transition inequalities in
Lemma~\ref{lem:uniform} uses a lower sum for $\delta(r_0)$.
Put $\ell'=1414213/10^6<\sqrt2<m'=1414214/10^6$,
$H_0=r_0/4096$, $r_j=jH_0$, and $b_j=\ell' r_j+r_j^2$.
As $\delta'$ is increasing and
$\delta'(r_j)\ge b_j/[m'(1+b_j)]$, one has
\[
 \delta(r_0)\ge \frac{H_0}{10^{12}}
 \sum_{j=0}^{4095}\left\lfloor\frac{10^{12} b_j}{m'(1+b_j)}\right\rfloor
 =\frac{1592131749274149}{40960000000000000}.
\]
This lower bound satisfies both comparisons at $r_0$.

\section*{Declaration of AI-assisted tools}
Author used ChatGPT solely to improve the readability and language of the manuscript. The final text was thoroughly reviewed by the author and take the full responsibility.

\begin{thebibliography}{99}
\addcontentsline{toc}{section}{\refname}
\small\raggedright
\setlength{\itemsep}{4pt plus 1pt minus 1pt}
\bibitem{HuKim2024} Hu, Su, Kim, Min-Soo (2024). Asymptotic expansions for the alternating Hurwitz zeta function and its derivatives. Journal of Mathematical Analysis and Applications, 537(1), 128306, \href{https://doi.org/10.1016/j.jmaa.2024.128306}{doi: 10.1016/j.jmaa.2024.128306}.

\bibitem{BalanzarioSanchez2012} Balanzario, Eugenio P., Sánchez-Ortiz, Jorge (2012). Riemann-Siegel integral formula for the Lerch zeta function. Mathematics of Computation, 81(280), 2319--2333, \href{https://doi.org/10.1090/S0025-5718-2011-02566-4}{doi: 10.1090/S0025-5718-2011-02566-4}.

\bibitem{Balanzario2004} Balanzario, Eugenio P. (2004). A Riemann-Siegel formula for the Hurwitz zeta function. Boletín de la Sociedad Matemática Mexicana, Serie 3, 10(1), 1--13, \url{https://www.smm.org.mx/boletin_anterior/v10/n1.pdf}.

\bibitem{GarunkstisLaurincikasSteuding2003} Garunkštis, Ramūnas, Laurinčikas, Antanas, Steuding, Jörn (2003). An approximate functional equation for the Lerch zeta function. Mathematical Notes, 74(4), 469--476, \href{https://doi.org/10.1023/A:1026183524830}{doi: 10.1023/A:1026183524830}.

\bibitem{Miyagawa2017} Miyagawa, Takashi (2017). Approximate functional equations for the Hurwitz zeta-function and Lerch zeta-function. Commentarii Mathematici Universitatis Sancti Pauli, 66(1-2), 15--27, \url{https://arxiv.org/abs/1704.01850}.

\bibitem{NavasPalencia2023} Navas-Palencia, Guillermo (2023). Numerical methods and arbitrary-precision computation of the Lerch transcendent. \url{https://arxiv.org/abs/2302.05928}.

\bibitem{FerreiraLopez2004} Ferreira, Chelo, López, José L. (2004). Asymptotic expansions of the Hurwitz--Lerch zeta function. Journal of Mathematical Analysis and Applications, 298(1), 210--224, \href{https://doi.org/10.1016/j.jmaa.2004.05.040}{doi: 10.1016/j.jmaa.2004.05.040}.

\bibitem{FernandezFokas2018} Fernandez, Arran, Fokas, Athanassios S. (2018). Asymptotics to all orders of the Hurwitz zeta function. Journal of Mathematical Analysis and Applications, 465(1), 423--458, \href{https://doi.org/10.1016/j.jmaa.2018.05.012}{doi: 10.1016/j.jmaa.2018.05.012}.

\bibitem{FernandezSpenceFokas2018} Fernandez, Arran, Spence, E. A., Fokas, Athanassios S. (2018). Uniform asymptotics as a stationary point approaches an endpoint. IMA Journal of Applied Mathematics, 83(1), 204--242, \href{https://doi.org/10.1093/imamat/hxx042}{doi: 10.1093/imamat/hxx042}.

\bibitem{Nemes2017} Nemes, Gergő (2017). Error bounds for the asymptotic expansion of the Hurwitz zeta function. Proceedings of the Royal Society A, 473(2203), 20170363, \href{https://doi.org/10.1098/rspa.2017.0363}{doi: 10.1098/rspa.2017.0363}.

\bibitem{GuthmannSquare1997} Guthmann, Andreas (1997). New integral representations for the square of the Riemann zeta-function. Acta Arithmetica, 82(4), 309--330, \href{https://doi.org/10.4064/aa-82-4-309-330}{doi: 10.4064/aa-82-4-309-330}.

\bibitem{GuthmannCusp1997} Guthmann, Andreas (1997). The Riemann-Siegel integral formula for Dirichlet series associated to cusp forms. Analytic and Elementary Number Theory (Vienna, 1996), 53--69.

\bibitem{Paris2016} Paris, R. B. (2016). The Stokes phenomenon and the Lerch zeta function. Mathematica Aeterna, 6(2), 165--179, \url{https://arxiv.org/abs/1602.00099}.

\bibitem{CaiLopez2019} Cai, Xing Shi, López, José L. (2019). A note on the asymptotic expansion of the Lerch's transcendent. Integral Transforms and Special Functions, 30(10), 844--855, \href{https://doi.org/10.1080/10652469.2019.1627530}{doi: 10.1080/10652469.2019.1627530}.

\bibitem{Balanzario2006} Balanzario, Eugenio P. (2006). A generalized Euler-Maclaurin formula for the Hurwitz zeta function. Mathematica Slovaca, 56(3), 307--316, \url{https://dml.cz/manakin/handle/10338.dmlcz/131453}.

\bibitem{ButzerEtAl2011} Butzer, P. L., Ferreira, P. J. S. G., Schmeisser, G., Stens, R. L. (2011). The summation formulae of Euler--Maclaurin, Abel--Plana, Poisson, and their interconnections with the approximate sampling formula of signal analysis. Results in Mathematics, 59, 359--400, \href{https://doi.org/10.1007/s00025-010-0083-8}{doi: 10.1007/s00025-010-0083-8}.

\bibitem{Bleistein1966} Bleistein, Norman (1966). Uniform asymptotic expansions of integrals with stationary point near algebraic singularity. Communications on Pure and Applied Mathematics, 19(4), 353--370, \href{https://doi.org/10.1002/cpa.3160190403}{doi: 10.1002/cpa.3160190403}.

\bibitem{Apostol1951} Apostol, Tom M. (1951). On the Lerch zeta function. Pacific Journal of Mathematics, 1(2), 161--167, \url{https://msp.org/pjm/1951/1-2/pjm-v1-n2-s.pdf}.

\bibitem{ParisHurwitz2005} Paris, R. B. (2005). The Stokes phenomenon associated with the Hurwitz zeta function. Proceedings of the Royal Society A, 461, 297--304.

\bibitem{Riemann1859} Riemann, Bernhard (1859). Ueber die Anzahl der Primzahlen unter einer gegebenen Grösse. Monatsberichte der Berliner Akademie, 671--680, \url{https://www.claymath.org/collections/riemanns-1859-manuscript/}.

\bibitem{Siegel1932} Siegel, Carl Ludwig (1932). Über Riemanns Nachlass zur analytischen Zahlentheorie. Quellen und Studien zur Geschichte der Mathematik, Astronomie und Physik, 2, 45--80.

\bibitem{Titchmarsh1986} Titchmarsh, E. C. (1986). The Theory of the Riemann Zeta-Function. Oxford University Press.

\bibitem{Edwards1974} Edwards, Harold M. (1974). Riemann's Zeta Function. Academic Press.

\bibitem{Olver1997} Olver, Frank W. J. (1997). Asymptotics and Special Functions. A K Peters.

\bibitem{DLMF2026} NIST (2026). Digital Library of Mathematical Functions. \url{https://dlmf.nist.gov/2.4}, \url{https://dlmf.nist.gov/7.2}, \url{https://dlmf.nist.gov/25.11}, \url{https://dlmf.nist.gov/25.14}.

\bibitem{WolframMathematica} Wolfram Research, Inc. (2026). Mathematica, version 15.0.1. Champaign, IL. \url{https://www.wolfram.com/mathematica}.

\bibitem{SymPy2017} Meurer, Aaron, Smith, Christopher P., Paprocki, Mateusz, et al. (2017). SymPy: symbolic computing in Python. PeerJ Computer Science, 3, e103, \href{https://doi.org/10.7717/peerj-cs.103}{doi: 10.7717/peerj-cs.103}.
\end{thebibliography}

\medskip
{\small\noindent\textbf{Code and data.} The \href{https://github.com/tiagocampinasilva/alternating-hurwitz-zeta/blob/v1.0.0/Reproduce.nb}{Mathematica notebook}, data and verification scripts accompany this article in the \href{https://github.com/tiagocampinasilva/alternating-hurwitz-zeta}{GitHub repository}. The DOI for the computational supplement is \href{https://doi.org/10.5281/zenodo.24483843}{10.5281/zenodo.24483843}.\par}

\medskip
{\small\noindent\textit{Email:} \href{mailto:tiagocampinassilva@gmail.com}{tiagocampinassilva@gmail.com}\par}
\end{document}